% Définition du répertoire contenant les images
\graphicspath{{IMAGE/}}

% FRAME Intro
\begin{frame}

\includegraphics[width=12cm]{Logos.pdf}

\vfill

\begin{center}

\vspace*{1.5cm}

\LARGE
\textbf{Simulation}

\vspace*{1.5cm}
 DELHI GIS-R School


\large
9-12\up{th} April 2019

\vspace*{1.5cm}


\textbf{Hadrien Commenges \& Paul Chapron}

{\small

\vspace*{0.1cm}

\url{hadrien.commenges@univ-paris1.fr}

\url{paul.chapron@ign.fr}

}

\end{center}

\end{frame}

% FRAME
\begin{frame}{Stochastic simulation }

Stochastic simulation : data generation using \textbf{randomly drawn values}.

\begin{itemize}
\item \textbf{Monte Carlo:} generic term refering to process involving random process repetition.
\item \textbf{Bootstrap:} re-sampling methods (usually to estimate distribution) 
\item \textbf{Permutation:} reordering elements of a set   
\end{itemize}

~

\textbf{Why ? }

Most of the time:  to approach a distribution 

\begin{itemize}
\item (often) because the analytical way is hard 
\item because (sometimes) there is no analytical way 
\item because we look for robust estimation adequate to the use case data 
\end{itemize}

\end{frame}


% FRAME
\begin{frame}{Monte Carlo}

$\rightarrow$ \textbf{Example:} iterated dice rolls \\
$\rightarrow$ \textbf{Goal:} examplify the Law of Large Numbers

~

\textbf{Expected value $\mu$ of a dice roll:}

$$
\mu = \frac{1}{6} \times 1 + \frac{1}{6} \times 2 + \frac{1}{6} \times 3 + \frac{1}{6} \times 4 + \frac{1}{6} \times 5 + \frac{1}{6} \times 6
$$


$$
\mu = \frac{1+2+3+4+5+6}{6} = 3,5
$$

~

\textbf{(weak) Law of Large Numbers: }

~
\begin{center}
$\bar{X_n} \to \mu$ when $n \to \infty$
\end{center}

~

\emph{(sample average converges toward the expected value, for a sufficiently large sample)}

\end{frame}


% FRAME
\begin{frame}{Monte Carlo}
%TODO refaire figure EN
\begin{figure}
\includegraphics[width=11cm]{TirageDe.pdf}
\end{figure}

\end{frame}



% FRAME
\begin{frame}{Bootstrap}

\textbf{Classical inference:}
\begin{itemize}
  \item Goal : approach  $\mu$ and $\sigma$  \textbf{parameters} of a distribution
  \item $\bar{X}$ and $\sigma_X$ are computed on a \textbf{sample} X
  \item Sampling of X and $\bar{X}$ / $\sigma_X$ computation are repeated 
\end{itemize}

~

\textbf{Example:}
\begin{itemize}
  \item For a 12M population (people from Île-de-France).
  \item People travel daily  between  0 and 200 km.
  \item 500 samples are drawn,  100 people each.
  \item For each sample, the average travel distance ($\bar{X}$) is computed .
\end{itemize}

~

$\rightarrow$ these 500 average values form a \textit{distribution}: the \textbf{sampling distribution of the mean}

\end{frame}



% FRAME
\begin{frame}{Bootstrap}

\textbf{Sampling distribution of the mean} (500 mean values):
 
\begin{figure}
\includegraphics[width=11cm]{DensiteNormal_EN.pdf}
\end{figure}

where $\mu$ et $\sigma$ are the \textbf{parameters} --  real mean and standard deviation of the population -- and $n$ is the  \textbf{sample size}. 

~

$\rightarrow$ This is \textbf{central limit theorem} (\textbf{CLT}).

\end{frame}




% FRAME
\begin{frame}{Bootstrap}

\begin{itemize}
\item \textbf{General inference idea:} the sample (which is known) allows to approach the parameters of the population distribution (which is unknown)
\item \textbf{General bootstrap idea:} re-sampling (which is known) from the sample (which is known) give insights about what would sampling look like on the whole population.
\item \textbf{Example:} to estimate the variance of the sampling distribution , we compute the variance of the re-sampling distribution of the mean.
\end{itemize}

\end{frame}



% FRAME
\begin{frame}{Bootstrap}

\textbf{Daily travel distance for Île-de-France people}
%TODO figure EN
\begin{figure}
\includegraphics[width=11cm]{PopDist.png}
\end{figure}

\footnotesize
Source: Enquête Globale Transport 2010 (French transportation national census)
\normalsize

\end{frame}




% FRAME
\begin{frame}{Bootstrap}

\textbf{Sampling distribution of the mean}
%TODO figure EN

\begin{figure}
\includegraphics[width=11cm]{SampDist.png}
\end{figure}

\end{frame}


% FRAME
\begin{frame}{Bootstrap}

\textbf{Re-sampling of the mean} on a sample 

\begin{figure}
\includegraphics[width=11cm]{SampDistBoot.png}
\end{figure}

\end{frame}


% FRAME
\begin{frame}{Bootstrap}


\textbf{Power of the Bootstrap}: this technique offers

\begin{itemize}
\item compute estimates (mean, variance) without any hypothesis or \textit{a priori} knowledge on the population
\item compute estimates variability  (confidence interval) without any hypothesis or \textit{a priori} knowledge on the population (\textit{distribution-free confidence intervals})
\item assess the stability of some model results  (\textit{cross-validation})
\end{itemize}

\end{frame}



